% demo-deck.tex
%
% Prompt 015 demonstration deck. Exercises the CS2 slide system's major
% pattern types on a small, deliberately generic example (a Stack ADT)
% so the system can be validated without prematurely authoring real
% course content. This is NOT Week 3, and NOT Welcome to the Odyssey --
% see sidecar/prompts/015_build_cs2_slide_system.md, "Explicitly out of
% scope."

\documentclass[aspectratio=169]{beamer}
\usepackage{../theme/cs2theme}
\usepackage{../theme/cs2commands}

\setcsweek{Demo Deck -- Prompt 015}

\title{The CS2 Presentation System}
\subtitle{A demonstration deck exercising the reusable slide grammar}
\author{Computer Science II \textbullet\ Dr. Jeremy P. Evert}
\date{Fall 2026 \textbullet\ Prompt 015 validation}

\begin{document}

% Pattern 1: Title / opening
\begin{frame}[plain]
  \titlepage
\end{frame}

% Pattern 2: Today's destination
\begin{destinationframe}{Today's Destination}
  \begin{itemize}
    \item Explain what a \textbf{Stack} guarantees, in your own words
    \item Trace \texttt{push}/\texttt{pop} against a real call sequence
    \item Predict, then verify, the output of a buggy stack-based function
    \item Connect the Stack contract to this week's Odyssey checkpoint
  \end{itemize}
  \instructorcue{Write the four bullets on the board before slides go up -- it's the map for the whole class meeting.}
\end{destinationframe}

% Pattern 3: Why this matters
\begin{whymattersframe}{Why This Matters}
  Every undo button, every browser back-arrow, every call stack you've
  ever silently trusted is a Stack.

  \vspace{3mm}
  \textbf{Today's claim:} once you can \emph{prove} a data structure's
  contract instead of just using it, you can build your own -- and know
  when someone else's is lying to you.
\end{whymattersframe}

% Pattern 4: Core concept
\begin{conceptframe}{Core Concept: LIFO}
  A \textbf{Stack} is a collection with exactly one required guarantee:

  \begin{center}
    \Large\color{cs2primary}\textbf{Last In, First Out}
  \end{center}

  \vspace{2mm}
  Everything else -- array-backed, linked-list-backed, bounded,
  unbounded -- is an \emph{implementation detail} the contract does not
  promise you.
\end{conceptframe}

% Pattern 5: Concept with diagram
\begin{frame}{Concept + Diagram: push/pop}
\conceptdiagramkicker
\begin{conceptdiagrambody}
  \begin{concepttext}
    A Stack exposes two required operations:
    \begin{itemize}
      \item \texttt{push(item)} -- add to the top
      \item \texttt{pop()} -- remove and return the top
    \end{itemize}
    \vspace{2mm}
    Both are $O(1)$ -- if yours isn't, check your backing structure.
  \end{concepttext}
  \begin{conceptdiagram}
    \begin{tikzpicture}[node distance=2mm]
      \node[cs2class] (c) {"c"};
      \node[cs2class, above=of c] (b) {"b"};
      \node[cs2class, above=of b] (a) {"a"};
      \node[left=4mm of a] (top) {\small\color{cs2accent}\textbf{top}};
      \draw[cs2dataflow] (top) -- (a);
    \end{tikzpicture}
  \end{conceptdiagram}
\end{conceptdiagrambody}
\end{frame}

% Pattern 6: Code example
\begin{frame}[fragile]{Code Example: A Minimal Stack}
\codeframekicker
\begin{lstlisting}
class Stack:
    def __init__(self):
        self._items = []

    def push(self, item):
        self._items.append(item)

    def pop(self):
        return self._items.pop()
\end{lstlisting}
\end{frame}

% Pattern 7: Code walkthrough (progressive reveal via overlays)
\begin{frame}[fragile,t]{Code Walkthrough: Where's the Bug?}
\codewalkthroughkicker
\begin{onlyenv}<1>
\begin{lstlisting}
def undo_last_two(history):
    history.pop()
    history.pop()
    return history
\end{lstlisting}
\end{onlyenv}
\begin{onlyenv}<2>
\begin{lstlisting}
def undo_last_two(history):
    history.pop()   # what if len(history) == 1?
    history.pop()
    return history
\end{lstlisting}
\end{onlyenv}
\begin{onlyenv}<3>
  \instructorcue{Let a student say the word "exception" before you confirm it. Don't just tell them.}
  \vspace{2mm}
  \textbf{Question:} what does \texttt{pop()} do on an empty Stack, and
  who decided that?
\end{onlyenv}
\end{frame}

% Pattern 8: Predict before running
\begin{frame}[fragile]{Predict Before Running}
\predictkicker
\begin{lstlisting}
s = Stack()
s.push(1)
s.push(2)
s.push(3)
print(s.pop())
print(s.pop())
\end{lstlisting}
  \vspace{2mm}
  \textbf{Predict:} what prints? Write it down before we run it.
\end{frame}

% Pattern 9: Live coding prompt
\begin{livecodingframe}{Live Coding: Balanced Parentheses}
  Together, build a function that uses a Stack to check whether
  \texttt{"(a(b)c)"} has balanced parentheses and \texttt{"(a(b)c"} does not.

  \instructorcue{Type it live. Let the first attempt be wrong.}
\end{livecodingframe}

% Pattern 10: Worked problem
\begin{workedproblemframe}{Worked Problem: Reverse a String with a Stack}
  \begin{enumerate}
    \item Push every character of \texttt{"hello"} onto a Stack
    \item Pop every character off, appending to a new string
    \item What order does the new string come out in, and why?
  \end{enumerate}
\end{workedproblemframe}

% Pattern 11: Think / pair / discuss
\begin{pairdiscussframe}{Think \textbullet\ Pair \textbullet\ Discuss}
  With your partner: could you implement a Queue using two Stacks?
  \stopandwork{3 minutes. Sketch it on paper -- no code yet.}
\end{pairdiscussframe}

% Pattern 12: Common failure or misconception
\begin{misconceptionframe}{Common Failure: "Stacks Are Just Lists"}
  A Python \texttt{list} \emph{can implement} a Stack, but a Stack is not
  a list.

  \vspace{2mm}
  If your code calls \texttt{.insert(0, x)} or indexes into the middle,
  you have stopped using the Stack contract -- even if the underlying
  object is still a \texttt{list}.
\end{misconceptionframe}

% Pattern 13: Compare two approaches
\begin{frame}[fragile]{Compare: Array-Backed vs. Linked Stack}
\begin{comparebody}
  \begin{compareleft}{Array-backed (Python list)}
\begin{lstlisting}
class Stack:
    def __init__(self):
        self._d = []
    def push(self, x):
        self._d.append(x)
    def pop(self):
        return self._d.pop()
\end{lstlisting}
  \end{compareleft}
  \begin{compareright}{Linked-node}
\begin{lstlisting}
class Node:
    def __init__(self, v, nxt):
        self.v, self.nxt = v, nxt

class Stack:
    def __init__(self):
        self._top = None
    def push(self, x):
        self._top = Node(x, self._top)
\end{lstlisting}
  \end{compareright}
\end{comparebody}
\end{frame}

% Pattern 14: Design decision
\begin{designdecisionframe}{Design Decision: What Should pop() Do When Empty?}
  Three real options -- and real libraries disagree:
  \begin{itemize}
    \item Raise an exception (Python's own \texttt{list.pop()})
    \item Return a sentinel value like \texttt{None}
    \item Block/wait for an item (some concurrent queue designs)
  \end{itemize}
  \vspace{2mm}
  \textbf{Your World Bible must state which one your Stack chose, and why.}
\end{designdecisionframe}

% Pattern 15: AI collaboration / verification moment
\begin{aimomentframe}{AI Collaboration Moment}
  If you ask an AI tool to implement \texttt{Stack}, you still owe the
  same chain every time:

  \begin{center}
    \small baseline $\rightarrow$ bounded proposal $\rightarrow$ diff
    $\rightarrow$ independent test $\rightarrow$ read/reason
    $\rightarrow$ accept or reject
  \end{center}

  \vspace{2mm}
  An AI-generated \texttt{pop()} that silently returns \texttt{None} on
  empty is a design decision it made \emph{for you} -- catch it.
\end{aimomentframe}

% Pattern 16: Reasoning Odyssey connection
\begin{odysseyframe}{Reasoning Odyssey Connection}
  Today's Stack work feeds directly into this week's Odyssey capability:
  your software world needs something that remembers order and can undo
  itself.

  \vspace{2mm}
  \textbf{Capability unlocked:} bounded, ordered, reversible state.
\end{odysseyframe}

% Pattern 17: World Bible entry
\begin{worldbibleframe}{World Bible Entry}
  \begin{worldbibleentry}{Rule: The Undo Stack Never Silently Drops State}
    \wbfield{Assumption}{Every reversible action pushes exactly one entry.}
    \wbfield{Invariant}{\texttt{len(undo\_stack) == number of applied actions}}
    \wbfield{Revision}{2026-08-17: changed pop-on-empty from \texttt{None} to raising, after a paired-programming bug report.}
  \end{worldbibleentry}
\end{worldbibleframe}

% Pattern 18: Odyssey checkpoint / gate
\begin{odysseygateframe}{Odyssey Checkpoint: Ordered State}
  \textbf{Gate requirement:} your software world must contain at least
  one object whose behavior depends on \emph{order}, backed by a Stack
  or Queue you implemented and tested yourself.

  \vspace{2mm}
  \small See \texttt{rubrics/odyssey\_gates/} for the full evidence shape --
  this slide is the pointer, not the rubric.
\end{odysseygateframe}

% Pattern 19: Evidence students need to produce
\begin{evidenceframe}{Evidence to Produce}
  \begin{itemize}
    \item A tested \texttt{Stack} class with \texttt{push}/\texttt{pop}
    \item A World Bible entry naming your empty-pop design decision
    \item One paired-programming note: what your partner caught
  \end{itemize}
\end{evidenceframe}

% Pattern 20: Recap
\begin{recapframe}{Recap}
  \begin{itemize}
    \item A Stack's only promise is LIFO -- implementation is a choice
    \item \texttt{push}/\texttt{pop} should be $O(1)$
    \item Empty-pop behavior is a design decision, not an accident
    \item This feeds this week's ordered-state Odyssey checkpoint
  \end{itemize}
\end{recapframe}

% Pattern 21: Exit question / next move
\begin{exitframe}{Exit Question}
  \Large\color{cs2primary}
  \textbf{Could you explain to a partner, right now, why \texttt{list.pop(0)} is not the same promise as \texttt{list.pop()}?}

  \vspace{4mm}
  \normalsize\normalfont\color{cs2ink}
  \textbf{Next move:} bring a paired sketch of your Stack's World Bible
  entry to Wednesday.
\end{exitframe}

\end{document}
